% ======================================================================
%  DISCUSSION
% ======================================================================
\section{Discussion}\label{sec:discussion}

\subsection{Forecast Quality vs Economic Settlement}\label{sec:proxy_gap}

The paper's primary claim is forecast quality under a selective confidence gate: on the test set and the validation campaign, LightGBM improves directional accuracy and mean z-unit settlement over capacity-matched mechanical persistence (Section~\ref{sec:results}). The settlement proxy $\text{pnl}_{\text{proxy}} = \text{direction} \times z_{t+H}$ is designed for that claim---it scores whether the standardized spread level moves as predicted---not as a dollar P\&L.

When the same validation $\tau{=}0.9$ ranked-capacity trades are re-scored in basis points of realized spread change, policy ordering can invert (Table~\ref{tab:bps}): z-unit wins for the learned policy do not automatically translate into positive mean gross bps. That gap is expected under the training objective (predict $z_{t+H}$, which is dominated by the slow-moving rolling mean) rather than instantaneous $\Delta S_t$. We therefore report z-unit DirAcc, R$^2$, and mean $\texttt{pnl\_proxy}$ as the headline evidence, and treat fee-aware bps as a separate economic layer for future work.

\begin{table}[h]
\caption{Validation $\tau{=}0.9$ ranked-capacity entries scored in z-units and basis points of realized spread change, gross of fees. Units answer different questions: z-proxy for forecast quality, bps for raw spread change.}
\label{tab:bps}
\begin{tabular}{lrrr}
\toprule
\textbf{Policy} & \textbf{Mean} & \textbf{Mean gross} & \textbf{Gross} \\
 & \textbf{z-proxy} & \textbf{(bps)} & \textbf{win rate} \\
\midrule
LightGBM              & $+1.195$ & $-0.94$ & 34.9\% \\
Mech.\ persistence    & $+0.731$ & $-3.97$ & 9.9\% \\
Mech.\ mean-reversion & $-0.731$ & $+3.97$ & 76.0\% \\
\bottomrule
\end{tabular}
\end{table}

\subsection{Portfolio Stability Complements Per-Trade Skill}

Under the $\tau{=}0.9$ ranked-capacity baseline book, capacity-matched mechanical persistence posts a higher hourly portfolio Sharpe (2.62) than LightGBM (1.90). That ranking does \emph{not} overturn the model's advantage. Hourly Sharpe scores the path of \emph{aggregated} hourly z-proxy, not forecast quality per fill: persistence admits about $4{\times}$ as many closes ($n{\approx}8{,}550$ vs $2{,}026$) because $|z_t|{\geq}0.9$ fires more often than $|\hat{z}|{\geq}0.9$, so its hourly P\&L series is thicker and can look stabler even though each trade is weaker. On the metrics that match the paper claim---per-trade DirAcc (84.9\% vs 69.5\%) and mean z-proxy ($+1.20$ vs $+0.73$) under the same gate and capacity---LightGBM remains clearly ahead. We therefore treat hourly Sharpe as a secondary stability check and rank policies by selective forecast skill.

